\pdfinfoomitdate=1
\pdftrailerid{}
\pdfsuppressptexinfo=15
\documentclass[11pt,letterpaper]{article}
\usepackage[T1]{fontenc}
\usepackage{lmodern,microtype}
\usepackage[margin=1in]{geometry}
\usepackage{amsmath,amssymb,amsthm,mathtools}
\usepackage{booktabs,longtable,array}
\usepackage[dvipsnames]{xcolor}
\usepackage{enumitem}
\usepackage[colorlinks=true,linkcolor=MidnightBlue,urlcolor=MidnightBlue,citecolor=MidnightBlue,bookmarksnumbered=true]{hyperref}
\usepackage{fancyhdr}
\pdfmapfile{+lm.map}\pdfmapfile{+cm.map}\pdfmapfile{+symbols.map}
\hypersetup{pdftitle={Report 230 Corrected Poisson Endpoint Expansions},pdfauthor={Research report},pdfsubject={Fixed order counts, marked laws, and controlled inverses}}
\setlength{\headheight}{14pt}
\pagestyle{fancy}\fancyhf{}
\fancyhead[L]{\small Poisson endpoint expansions}
\fancyhead[R]{\small Report 230}
\fancyfoot[C]{\thepage}
\setlength{\parskip}{4pt}
\setlength{\emergencystretch}{2em}
\linespread{1.06}
\numberwithin{equation}{section}
\newtheorem{theorem}{Theorem}[section]
\newtheorem{proposition}[theorem]{Proposition}
\newtheorem{lemma}[theorem]{Lemma}
\newtheorem{corollary}[theorem]{Corollary}
\theoremstyle{definition}\newtheorem{definition}[theorem]{Definition}
\theoremstyle{remark}\newtheorem{remark}[theorem]{Remark}
\newcommand{\E}{\mathbb E}
\newcommand{\Pp}{\mathbb P}
\newcommand{\Pois}{\operatorname{Pois}}
\newcommand{\Var}{\operatorname{Var}}
\newcommand{\Cov}{\operatorname{Cov}}
\newcommand{\TV}{d_{\mathrm{TV}}}
\newcommand{\1}{\mathbf 1}
\newcommand{\Rcal}{\mathcal R}
\newcommand{\N}{\mathbb N}
\newcommand{\Q}{\mathbb Q}
\newcommand{\ii}{\mathrm i}
\DeclareMathOperator{\Res}{Res}
\title{Corrected Poisson Endpoint Expansions\\[4pt]\large Critical resummation, inactive vertices, and controlled index inverses\\for a family of self powered binomial sums}
\author{Report 230}
\date{5 October 2026}
\begin{document}
\maketitle
\begin{abstract}
For the exact sums
\[
 a_p(n)=\sum_{0\le k\le n/(p+1)}(n-pk)^{pk}\binom{n-pk}{k},
\]
we prove every fixed algebraic order for each fixed integer $p\ge1$.
An exact residue-conditioned Poisson representation gives a small deformation
when $p\ge2$ and a critical, nonvanishing deformation when $p=1$.
The latter requires a centered resummation and has half-power corrections,
with no odd quarter powers. We calculate corrected coefficients for OEIS
A360592, A360479, and A360747; their observed leading equivalents remain valid.
For $p\ge2$ we prove weighted density expansions, mean and variance refinements,
and sharp total-variation constants for the inactive-vertex count, including an
improved tilted Poisson reference. Explicit finite smooth models yield controlled
Newton and convergent local Lagrange inverses, arbitrary fixed index accuracy,
and a two-ceiling envelope for exact thresholds. The proofs include global
tail bounds and never differentiate an unspecified discrete remainder.
A self-contained exact-arithmetic package accompanies the article.
\end{abstract}
\noindent\textbf{Scope.} All asymptotic assertions keep $p$ and the requested order fixed.
No secondary exponential sector, optimal truncation theorem, growing-$p$ limit,
effective finite inverse onset, or critical marked-law theorem is asserted.
\tableofcontents
\section{Results and reading guide}
\label{sec:overview}
The family under consideration is defined for integers $p\ge1$ and $n\ge0$ by
\begin{equation}\label{eq:counts}
 a_p(n)=\sum_{k=0}^{\lfloor n/q\rfloor}(n-pk)^{pk}\binom{n-pk}{k},
 \qquad q=p+1.
\end{equation}
We set the $k=0$ summand equal to one, including at $n=0$.
Thus $a_p(0)=1$; no normalization involving $n/q$ is used at zero.
The cases $p=1,2,3$ are the exact finite sums in A360592, A360479, and
A360747 respectively~\cite{oeis1,oeis2,oeis3}.

The key parameter is
\begin{equation}\label{eq:lambda}
 c_p=\exp(p^2/q)q^{-1/q},\qquad
 \lambda=\lambda_p(n)=c_p n^{1/q}.
\end{equation}
For $p\ge2$, let
\begin{equation}\label{eq:normalization}
 M_p(n)=\frac1q(n/q)^{pn/q}e^\lambda,
 \qquad \Rcal_p(n)=\frac{a_p(n)}{M_p(n)}.
\end{equation}
The first polynomial perturbation has a particularly simple universal form:
\begin{equation}\label{eq:firstuniv}
 \Rcal_p(n)=1-\frac{A\lambda^2+b\lambda}{n}
       +O_p\!\left(\frac{\lambda^4}{n^2}\right),\qquad
 A=\frac{p^2+1/q}{2},\quad b=A-\frac q2=\frac{p(p^2-2)}{2q}.
\end{equation}
For every fixed $K\ge0$, the stronger expansion is
\begin{equation}\label{eq:overviewfixed}
 \Rcal_p(n)=\sum_{j=0}^K C_{p,j}(c_p)n^{-j/q}
                  +O_{p,K}(n^{-(K+1)/q}),
\end{equation}
where $C_{p,j}$ is an explicitly generated rational polynomial and
$C_{p,0}=1$. The coefficients $C_{p,1},\ldots,C_{p,p-2}$ vanish.
Every assertion is uniform over the finitely many residue classes of $n$ modulo $q$.

At $p=1$ the quantity $\lambda^2/n$ is constant, so
\eqref{eq:firstuniv} is not an asymptotic expansion around one. The correct
normalization is instead
\begin{equation}\label{eq:criticalM}
 c=c_1=\sqrt{e/2},\quad d=3e/8=3c^2/4,\qquad
 M_1(n)=\frac12(n/2)^{n/2}\exp(c\sqrt n-d).
\end{equation}
Then $a_1(n)/M_1(n)$ has an expansion in $n^{-1/2}$ to every fixed
order. There are no odd quarter powers. Its first correction is
\begin{equation}\label{eq:criticalfirst}
 C_{1,1}(c)=\frac c4+\frac{11c^3}{8}.
\end{equation}

\subsection{What is corrected}
The observed OEIS formulas all use the symbol $\sim$. Interpreted literally
as a leading equivalent, each remains correct. Their displayed
parenthesized refinements suggest more precise coefficients; it is these
intended coefficients that disagree with the proved expansions here.
For $p=1$ and $p=2$ the observed first coefficient exceeds the correct one
by $1/(8c_p)$. For $p=3$ the observed coefficient of $n^{-1/4}$ is
$1/(8c_3)$, whereas the true coefficient is zero.
The first nonzero $p=3$ correction is negative and of order $n^{-1/2}$.
These statements concern inspected or indexed source snapshots, not
verified latest revision histories. Section~\ref{sec:sources} records the
source scope explicitly.

\subsection{What the extra results add}
The integer $r=n-qk$ has a useful combinatorial meaning: it counts inactive
vertices in a canonically labeled directed-graph model. For fixed $p\ge2$,
its distribution is close to a Poisson distribution conditioned on
$r\equiv n\pmod q$. We identify the sharp total-variation constants and
improve the approximation by an explicit Lambert-$W$ tilt. The conditioning
cannot be dropped in total variation.

The rapidly increasing counts also have controlled inverse descriptions.
An explicit finite asymptotic model, followed by finitely many Newton steps,
recovers the index on the exact range to any prescribed fixed algebraic
accuracy. For a general real level $y$, we prove an envelope involving two
ceilings. It can leave two adjacent possible thresholds; arbitrary tiny
real error does not justify an unqualified rounding rule.

\subsection{Proof architecture and classical ingredients}
Section~\ref{sec:exact} obtains an exact Poisson representation by elementary
algebra. Section~\ref{sec:subcritical} combines global domination, a local
Cauchy estimate, and the root-of-unity filter. Section~\ref{sec:critical}
uses a different centered argument at $p=1$.
Section~\ref{sec:marked} proves the marked laws with weighted remainders.
Section~\ref{sec:inverse} treats finite smooth inverses.

The underlying tools are classical: Poisson factorial and central moments,
finite-difference or Faulhaber polynomial calculus, root-of-unity filters,
Chernoff bounds, analytic Taylor estimates, Charlier polynomials,
Lagrange inversion, and the Lambert $W$ function~\cite{dlmf,corless}.
We give the identities and estimates needed here so that the argument does
not depend on locating a theorem with this exact sequence as an example.
No first-priority claim is needed for any mathematical conclusion.

\section{Exact counting and Poisson reduction}
\label{sec:exact}
\subsection{A formal generating function and a graph model}
Expanding the finite binomial at each $m$ gives the formal identity
\begin{equation}\label{eq:ogf}
 \sum_{n\ge0}a_p(n)x^n
       =\sum_{m\ge0}x^m\bigl(1+(mx)^p\bigr)^m.
\end{equation}
Indeed, the term indexed by $k$ on the right is
$\binom mk m^{pk}x^{m+pk}$. Setting $m=n-pk$ gives
\eqref{eq:counts}, and $k\le m$ becomes $qk\le n$.
Every coefficient receives only finitely many contributions, so no analytic
convergence is being claimed. In fact, the coefficients have superexponential
growth.

For a direct interpretation, use the canonical labeled vertex set
$[m]=\{1,\ldots,m\}$. Choose $k$ active vertices. Each active vertex has $p$
ordered, or color-labeled, outgoing arc slots, and each slot independently
chooses a target in $[m]$. Loops and repeated targets are allowed.
Inactive vertices have no outgoing slots. There are
$\binom mk m^{pk}$ such structures. The size statistic is
\[
 n=m+pk=\hbox{number of vertices plus number of arcs}.
\]
The disjoint union over the possible $m$ therefore has size $a_p(n)$.
This is a labeled model with a nonstandard size statistic;
\eqref{eq:ogf} is not an ordinary unlabeled-digraph generating function.

The number of inactive vertices is
\begin{equation}\label{eq:rdef}
 r=m-k=n-qk,\qquad
 k=\frac{n-r}{q},\qquad m=\frac{n+pr}{q}.
\end{equation}
For a uniformly chosen structure of size $n$, denote this random number by
$R_n$. Its support is exactly
$\{r:0\le r\le n,\ r\equiv n\pmod q\}$.

\subsection{An exact change of variables}
For $n>0$ write $N=n/q$ and $B=N^{pn/q}$.
For an allowed $r$, use $\binom mk=\binom mr$ to get
\[
 m^{pk}\binom mk
  =\frac{m^{pk+r}}{r!}\prod_{h=0}^{r-1}\left(1-\frac hm\right).
\]
Since $pk+r=(pn+r)/q$ and $m=N(1+pr/n)$, this becomes
\begin{equation}\label{eq:exactterm}
 m^{pk}\binom mk=B\frac{\lambda^r}{r!}e^{\Delta_n(r)},
\end{equation}
where
\begin{equation}\label{eq:Delta}
 \Delta_n(r)=\frac{pn+r}{q}\log(1+pr/n)-\frac{p^2r}{q}
      +\sum_{h=0}^{r-1}\log\left(1-\frac{qh}{n+pr}\right).
\end{equation}
All factors inside the logarithms are positive for integers $0\le r\le n$:
$n+pr-q(r-1)=n-r+q>0$.
The empty product and empty sum at $r=0$ are one and zero respectively.
The same expression defines a real function at every integer in this range,
not only those satisfying the congruence.

If $X\sim\Pois(\lambda)$, summation gives the exact identity
\begin{equation}\label{eq:exactpoisson}
 q e^{-\lambda}\frac{a_p(n)}{B}
  =q\E\left[e^{\Delta_n(X)}\1_{\{X\le n\}}
                    \1_{\{X\equiv n\ (q)\}}\right].
\end{equation}
For $p\ge2$ the left side is $\Rcal_p(n)$.
For $p=1$ it is $e^{-d}a_1(n)/M_1(n)$.
The factor $q$ is the lattice spacing of the residue condition. It does
not come from a Stirling approximation; no approximation has yet been used.

\subsection{Poisson polynomial identities}
Let $T_j$ denote the Touchard polynomial,
\begin{equation}\label{eq:touchard}
 T_j(z)=\sum_{\ell=0}^j\left\{\begin{matrix}j\\\ell\end{matrix}\right\}z^\ell,
 \qquad \E X^j=T_j(\lambda).
\end{equation}
Here the braces are Stirling numbers of the second kind.
This follows by writing $X^j$ in the falling-factorial basis and using
$\E(X)_\ell=\lambda^\ell$.
In particular, for each fixed $j$,
$\E X^j=O_j((1+\lambda)^j)$.

Put $\omega=e^{2\pi\ii/q}$ and
$\kappa_q=1-\cos(2\pi/q)>0$.
For any integer residue $a$ and any nonnegative integer $j$,
\begin{equation}\label{eq:filter}
 q\E[X^j\1_{\{X\equiv a\ (q)\}}]
  =\sum_{\ell=0}^{q-1}\omega^{-a\ell}
       e^{\lambda(\omega^\ell-1)}T_j(\lambda\omega^\ell).
\end{equation}
To see this, insert
$q\1_{\{X\equiv a\}}=\sum_{\ell=0}^{q-1}\omega^{\ell(X-a)}$
and use the Poisson generating function. The $\ell=0$ term is $T_j(\lambda)$;
all other terms are
$O_{q,j}((1+\lambda)^j e^{-\kappa_q\lambda})$.
For a fixed polynomial, residue restriction therefore changes its Poisson
average only by an exponentially small absolute error.
This statement concerns polynomial averages. It does not identify a
relative secondary exponential sector of the full nonpolynomial sum.

\subsection{Exact monotonicity}
\begin{proposition}\label{prop:monotone}
For every $p\ge1$,
\[
 a_p(0)=\cdots=a_p(p)=1,\qquad a_p(p+1)=2,
\]
and $a_p(n+1)>a_p(n)$ for every $n\ge p$.
Consequently, for each real $y>1$ the threshold
\begin{equation}\label{eq:threshold}
 T_p(y)=\min\{n\ge0:a_p(n)\ge y\}
\end{equation}
is well-defined.
\end{proposition}
\begin{proof}
For each previously allowed $k$, increasing $n$ increases
$m=n-pk\ge k$. Both $m^{pk}$ and $\binom mk$ are nondecreasing in $m$,
and all newly allowed summands are positive. At the first step $n=p$,
the new $k=1$ term gives strictness. Thereafter the common $k=1$ term
$(n-p)^{p+1}$ strictly increases. It also proves unboundedness.
The initial equalities are immediate from \eqref{eq:counts}.
\end{proof}

\section{Every fixed order for the small deformation regime}
\label{spb:ep:sec:subcritical}
Throughout this section $p\ge2$ is a fixed integer. Thus
$\lambda\asymp_p n^{1/q}$ and
\begin{equation}\label{spb:ep:eq:small}
 \lambda^2/n\asymp_p n^{-(p-1)/q}\longrightarrow0.
\end{equation}
This is the small parameter that makes an ordinary exponential expansion
of the defect valid after Poisson averaging.

\subsection{Global domination and tails}
\begin{lemma}\label{spb:ep:lem:sign}
For every integer $0\le r\le n$, $\Delta_n(r)\le0$.
Consequently
\begin{equation}\label{spb:ep:eq:globalR}
 0<\Rcal_p(n)\le q\Pp(X\equiv n\pmod q)
             =1+O_p(e^{-\kappa_q\lambda}).
\end{equation}
\end{lemma}
\begin{proof}
Let
\[
 F(r)=\frac{pn+r}{q}\log(1+pr/n)-\frac{p^2r}{q}.
\]
Direct differentiation gives $F(0)=F'(0)=0$ and
\begin{equation}\label{spb:ep:eq:Fsecond}
 F''(r)=\frac{p(2-p^2+pr/n)}{qn(1+pr/n)^2}.
\end{equation}
For $0\le r\le n$, its numerator is at most
$p(2-p^2+p)\le0$. Hence $F(r)\le0$.
Every logarithm in the sum in \eqref{spb:ep:eq:Delta} is nonpositive.
This proves the sign. Now apply \eqref{spb:ep:eq:exactpoisson} and
\eqref{spb:ep:eq:filter} with $j=0$.
\end{proof}

The right side of \eqref{spb:ep:eq:globalR} need not be identically one:
residue probabilities may have exponentially small oscillations.
Nevertheless it already excludes a positive algebraic first correction.
In particular, a positive $n^{-1/4}$ coefficient for $p=3$ is impossible.

The Poisson moment generating function gives, with
$h=2\log2-1>0$,
\begin{equation}\label{spb:ep:eq:tail}
 \Pp(X>2\lambda)\le e^{-h\lambda}.
\end{equation}
Indeed, exponential Markov with parameter $\log2$ bounds the probability
by $\exp(\lambda(2-1)-2\lambda\log2)$.
For each fixed integer $d\ge0$, Cauchy--Schwarz yields
\begin{equation}\label{spb:ep:eq:weightedtail}
 \E[X^d\1_{\{X>2\lambda\}}]
   \le (\E X^{2d})^{1/2}\Pp(X>2\lambda)^{1/2}
   =O_d((1+\lambda)^d e^{-h\lambda/2}).
\end{equation}
By Lemma~\ref{spb:ep:lem:sign}, the exact normalized contribution outside
$X\le2\lambda$ is bounded by the same unweighted tail, up to $q$.
For large $n$, $2\lambda<n$, so the original cutoff at $n$ causes no
additional difficulty. These are global estimates, including the far
endpoint $r=n$.

\subsection{Defect polynomials and a uniform analytic remainder}
For $j\ge1$ define
\begin{equation}\label{spb:ep:eq:Dgeneral}
 D_j(r)=(-1)^{j+1}\left\{
  \frac{p^j}{q}\left(\frac1j-\frac{p^2}{j+1}\right)r^{j+1}
  +\frac1j\sum_{h=0}^{r-1}\bigl[(pr-qh)^j-(pr)^j\bigr]\right\}.
\end{equation}
The finite sum is interpreted as its polynomial in $r$, obtained for
example from Faulhaber sums or interpolation in the falling-factorial
basis. Thus $D_j\in\Q[r]$, $D_j(0)=0$, and
$\deg D_j\le j+1$.
Define
\begin{equation}\label{spb:ep:eq:Hrec}
 H_0(r)=1,\qquad
 H_j(r)=\frac1j\sum_{s=1}^j sD_s(r)H_{j-s}(r).
\end{equation}
Formal differentiation of an exponential shows that
\[
 \exp\left(\sum_{j\ge1}D_j(r)z^j\right)
        =\sum_{j\ge0}H_j(r)z^j.
\]
Induction gives $\deg H_j\le2j$ and $H_j(0)=0$ for $j\ge1$.

\begin{lemma}\label{spb:ep:lem:cauchy}
For each fixed $S\ge0$, uniformly over integers $0\le r\le2\lambda$,
\begin{equation}\label{spb:ep:eq:pointrem}
 e^{\Delta_n(r)}=\sum_{j=0}^S H_j(r)n^{-j}
       +O_{p,S}\bigl((r+1)^{2S+2}n^{-S-1}\bigr).
\end{equation}
\end{lemma}
\begin{proof}
For an integer $r\ge1$ consider the analytic germ
\begin{equation}\label{spb:ep:eq:G}
 G_r(z)=e^{-p^2r/q}
       (1+prz)^{p/(qz)-pr/q}
       \prod_{h=0}^{r-1}(1+(pr-qh)z),
\end{equation}
using logarithmic branches at $z=0$. The apparent singularity at zero is
removable and $G_r(0)=1$. Recombining the denominators in
\eqref{spb:ep:eq:Delta} gives $G_r(1/n)=e^{\Delta_n(r)}$.

Choose a small constant $\eta_p>0$ and put
$\rho_r=\eta_p/(r+1)^2$. For $|z|\le\rho_r$, every logarithm argument in
\eqref{spb:ep:eq:G} lies in a fixed disk about one excluding zero. After its
constant $p^2r/q$ is subtracted, the expression
$(p/(qz))\log(1+prz)$ is $O_p(r^2|z|)$.
The remaining logarithm term and the sum of the $r$ product logarithms
are also $O_p(r^2|z|)$. Therefore
\[
 |\log G_r(z)|\le C_p r^2|z|\le C'_p,
 \qquad |G_r(z)|\le e^{C'_p},
\]
uniformly in $r$.
Expanding its logarithm gives exactly \eqref{spb:ep:eq:Dgeneral}; exponentiation
gives \eqref{spb:ep:eq:Hrec}.

Cauchy's coefficient estimate on $|z|=\rho_r$ and a geometric tail give
\[
 \left|G_r(1/n)-\sum_{j=0}^S H_j(r)n^{-j}\right|
  \le C_{p,S}(r+1)^{2S+2}n^{-S-1}
\]
whenever $(r+1)^2/n\le\eta_p/2$.
This condition holds uniformly for $r\le2\lambda$ by \eqref{spb:ep:eq:small}.
For $r=0$, $G_0=1$ and the conclusion is immediate.
\end{proof}

\subsection{The coefficient theorem}
Define the Poisson transform of $H_j$ by
\begin{equation}\label{spb:ep:eq:U}
 H_j(r)=\sum_d h_{j,d}r^d,
 \qquad U_j(z)=\sum_d h_{j,d}T_d(z).
\end{equation}
Thus $U_0=1$, $U_j(0)=0$ for $j\ge1$, and $\deg U_j\le2j$.

\begin{theorem}\label{spb:ep:thm:allp}
For every fixed $S\ge0$,
\begin{equation}\label{spb:ep:eq:allH}
 \Rcal_p(n)=\sum_{j=0}^S n^{-j}U_j(\lambda)
   +O_{p,S}\left(\frac{(1+\lambda)^{2S+2}}{n^{S+1}}
                              +e^{-c_{p,S}\lambda}\right)
\end{equation}
for some $c_{p,S}>0$. In particular the exponential term can be absorbed
into the algebraic remainder. Setting $t=n^{-1/q}$ and $\lambda=c_p/t$,
for every fixed $K\ge0$ this implies \eqref{spb:ep:eq:overviewfixed}, with
\begin{equation}\label{spb:ep:eq:Cextract}
 C_{p,k}(c)=[t^k]\sum_{j=0}^{\lfloor K/(p-1)\rfloor}
                         t^{qj}U_j(c/t),\qquad 0\le k\le K.
\end{equation}
\end{theorem}
\begin{proof}
Apply \eqref{spb:ep:eq:pointrem} in the exact expectation on $X\le2\lambda$.
The expectation of the remainder is at most a constant times
$n^{-S-1}\E(X+1)^{2S+2}$, giving the algebraic term in \eqref{spb:ep:eq:allH}.
The exact tail is controlled by \eqref{spb:ep:eq:tail} and the sign lemma.
The tails of the finitely many polynomials $H_j$ are controlled by
\eqref{spb:ep:eq:weightedtail}. Remove their residue restrictions using
\eqref{spb:ep:eq:filter}. All resulting exponential terms, including fixed
polynomial factors, can be combined into $O(e^{-c_{p,S}\lambda})$.

A monomial $n^{-j}\lambda^d$ becomes $c^dt^{qj-d}$, with
$1\le d\le2j$ when $j\ge1$. Its possible exponent lies in
\begin{equation}\label{spb:ep:eq:bands}
 [(p-1)j,\ qj-1].
\end{equation}
Take $S=\lfloor K/(p-1)\rfloor$.
The remainder in \eqref{spb:ep:eq:allH} starts at order
$t^{(p-1)(S+1)}=O(t^{K+1})$.
Discarding finitely many polynomial terms with exponent larger than $K$
also gives $O(t^{K+1})$. This proves \eqref{spb:ep:eq:Cextract} and the theorem.
\end{proof}

Formula \eqref{spb:ep:eq:bands} gives a structural exclusion of low powers.
It does not imply that every exponent inside an allowed band has a
nonzero coefficient. In particular, the absence of
$t,t^2,\ldots,t^{p-2}$ requires no symbolic computation.

\subsection{Explicit corrected terms}
The first two defect polynomials are
\begin{align}
 D_1(r)&=-Ar^2+\frac q2r,\label{spb:ep:eq:D1}\\
 D_2(r)&=\gamma_3r^3-\gamma_2r^2-\gamma_1r,\label{spb:ep:eq:D2}\\
 \gamma_3&=\frac{p^3}{3}-\frac1{6q},\qquad
 \gamma_2=\frac{p^2-1}{4},\qquad \gamma_1=\frac{q^2}{12}.\nonumber
\end{align}
Since $U_1(\lambda)=\E D_1(X)=-A\lambda^2-b\lambda$,
Theorem~\ref{spb:ep:thm:allp} gives \eqref{spb:ep:eq:firstuniv}.
For $p=2$, write $c=e^{4/3}/3^{1/3}$ and $t=n^{-1/3}$. Then
\begin{align}\label{spb:ep:eq:p2}
 \Rcal_2(n)={}&1-\frac{13c^2}{6}t
 +\left(\frac{169c^4}{72}-\frac{2c}{3}\right)t^2\nonumber\\
 &+\left(\frac{121c^3}{9}-\frac{2197c^6}{1296}\right)t^3
 +O(t^4).
\end{align}
One additional coefficient is
\[
 C_{2,4}(c)=\frac{134c^2}{9}-\frac{2977c^5}{108}
                         +\frac{28561c^8}{31104}.
\]
For $p=3$, write $c=e^{9/4}/\sqrt2$ and $t=n^{-1/4}$. Then
\begin{align}\label{spb:ep:eq:p3}
 \Rcal_3(n)={}&1-\frac{37c^2}{8}t^2-\frac{21c}{8}t^3
        +\frac{1369c^4}{128}t^4+O(t^5)\nonumber\\
 ={}&1-\frac{37e^{9/2}}{16}n^{-1/2}
       -\frac{21e^{9/4}}{8\sqrt2}n^{-3/4}
       +\frac{1369e^9}{512}n^{-1}+O(n^{-5/4}).
\end{align}
The next coefficients are
\[
 C_{3,5}(c)=\frac{12265c^3}{192},\qquad
 C_{3,6}(c)=\frac{9471c^2}{128}-\frac{50653c^6}{3072}.
\]
All these identities are rational-polynomial outputs of the finite
recurrence, not numerical coefficient fits.

\begin{remark}[Formal order and useful numerical accuracy]
The theorem is asymptotic for each fixed truncation. It does not say that
larger order improves every finite-$n$ approximation. The constants in
\eqref{spb:ep:eq:p3} are large. At $n=10000$, even the fractional-power
truncation through $t^{12}$ has relative error about $-0.662$.
Section~\ref{spb:ep:sec:numerics} gives the complete-sum diagnostic and contrasts
it with another legitimate finite truncation of \eqref{spb:ep:eq:allH}.
\end{remark}

\section{The critical half power expansion}
\label{spb:ep:sec:critical}
In this section $p=1$, $q=2$, $c=\sqrt{e/2}$,
$t=n^{-1/2}$, and $\lambda=c/t$.
The normalization $M_1$ is \eqref{spb:ep:eq:criticalM}.
Since $\lambda^2/n=c^2$, the leading quadratic defect persists:
$\Delta_n(X)$ is approximately $-3(tX)^2/4$, whose limiting value is
$-3c^2/4$. Expanding $e^{\Delta_n(X)}$ as a small perturbation of one would
miss an order-one factor. We retain that factor and expand around the
Poisson center instead.

\begin{theorem}\label{spb:ep:thm:critical}
For every fixed integer $K\ge0$, there are explicit polynomials
$C_{1,j}\in\Q[c]$ such that
\begin{equation}\label{spb:ep:eq:criticalall}
 \frac{a_1(n)}{M_1(n)}
     =\sum_{j=0}^K C_{1,j}(c)n^{-j/2}
                    +O_K(n^{-(K+1)/2}).
\end{equation}
Both parities have the same coefficients. There are no odd quarter powers
in the normalized expansion. The first four nonconstant coefficients are
\begin{align}
 C_{1,1}(c)&=\frac c4+\frac{11c^3}{8},\label{spb:ep:eq:critC1}\\
 C_{1,2}(c)&=\frac{31c^2}{32}-\frac{319c^4}{96}
                                      +\frac{121c^6}{128},\label{spb:ep:eq:critC2}\\
 C_{1,3}(c)&=-\frac{5c}{96}-\frac{739c^3}{128}
              +\frac{46517c^5}{3840}-\frac{7381c^7}{1536}
              +\frac{1331c^9}{3072},\label{spb:ep:eq:critC3}\\
 C_{1,4}(c)&=-\frac{227c^2}{128}+\frac{61923c^4}{2048}
              -\frac{771061c^6}{15360}+\frac{4140367c^8}{184320}
              \nonumber\\[-2pt]
          &\hspace{18mm}-\frac{41261c^{10}}{12288}
                         +\frac{14641c^{12}}{98304}.\label{spb:ep:eq:critC4}
\end{align}
\end{theorem}

\subsection{A separate global bound}
The exact Poisson identity is still available, now in the form
\begin{equation}\label{spb:ep:eq:criticalexact}
 \frac{a_1(n)}{M_1(n)}
  =2e^{3c^2/4}\E\left[e^{\Delta_n(X)}\1_{\{X\le n\}}
                               \1_{\{X\equiv n\ (2)\}}\right].
\end{equation}
The nonpositivity statement of Lemma~\ref{spb:ep:lem:sign} is false at $p=1$;
for example, $\Delta_n(1)>0$.
The required replacement is a uniform upper bound.

\begin{lemma}\label{spb:ep:lem:criticalbound}
For every integer $0\le r\le n$,
\begin{equation}\label{spb:ep:eq:criticalbound}
 \Delta_n(r)\le-\frac{r^2}{4n}+\frac r{2n}
              =\frac{1-(r-1)^2}{4n}\le\frac1{4n}.
\end{equation}
\end{lemma}
\begin{proof}
The first two terms of \eqref{spb:ep:eq:Delta} now form
$F(r)=((n+r)/2)\log(1+r/n)-r/2$.
One has $F(0)=F'(0)=0$ and
$F''(r)=1/(2n(1+r/n))\le1/(2n)$, so $F(r)\le r^2/(4n)$.
Since $\log(1-x)\le-x$ for $0\le x<1$,
\[
 \sum_{h=0}^{r-1}\log\left(1-\frac{2h}{n+r}\right)
       \le-\frac{r(r-1)}{n+r}\le-\frac{r(r-1)}{2n}.
\]
The final inequality uses $r(r-1)\ge0$, which holds at every allowed
integer. Adding the bounds proves \eqref{spb:ep:eq:criticalbound}.
\end{proof}

Let
\begin{equation}\label{spb:ep:eq:criticalbulk}
 \mathcal C_n=\{|X-\lambda|\le\lambda^{3/4}\}.
\end{equation}
The usual exponential Markov argument for the Poisson generating function
implies, for a positive absolute constant $b_0$ and large $\lambda$,
\begin{equation}\label{spb:ep:eq:criticaltail}
 \Pp(\mathcal C_n^c)
 \le2\exp\left(-\frac{\lambda^{3/2}}
                         {2(\lambda+\lambda^{3/4})}\right)
 =O(e^{-b_0\sqrt\lambda}).
\end{equation}
For completeness, the upper-tail Chernoff rate is
$\lambda((1+s)\log(1+s)-s)$ at displacement $s\lambda$;
integrating $\log(1+s)\ge s/(1+s)$ bounds it below by
$\lambda s^2/(2(1+s))$. The lower-tail rate is at least
$\lambda s^2/2$. Substitute $s=\lambda^{-1/4}$.

Lemma~\ref{spb:ep:lem:criticalbound} bounds the exact normalized contribution
on $\mathcal C_n^c$ by a fixed multiple of \eqref{spb:ep:eq:criticaltail}.
Every fixed polynomially weighted tail is also negligible by
Cauchy--Schwarz and Poisson moment bounds. Since $\lambda\asymp\sqrt n$,
these errors are smaller than every power of $n^{-1}$.
For large $n$, the central event lies inside $0\le X\le2\lambda<n$.

\subsection{Defect polynomials at criticality}
Specializing \eqref{spb:ep:eq:Dgeneral} gives
\begin{equation}\label{spb:ep:eq:Dcrit}
 D_j(r)=(-1)^{j+1}\left\{
  \frac{r^{j+1}}{2j(j+1)}
  +\frac1j\sum_{h=0}^{r-1}\bigl[(r-2h)^j-r^j\bigr]\right\}.
\end{equation}
The first values are
\begin{align*}
 D_1&=-3r^2/4+r,& D_2&=r^3/4-r/3,\\
 D_3&=-7r^4/24+r^3/3,&
 D_4&=7r^5/40-r^3/3+2r/15,\\
 D_5&=-11r^6/60+r^5/5.
\end{align*}
We require a local logarithmic estimate on a larger analytic scale than
in Lemma~\ref{spb:ep:lem:cauchy}.
For every fixed $L\ge1$, uniformly for integers $0\le r\le2\lambda$,
\begin{equation}\label{spb:ep:eq:criticalDrem}
 \Delta_n(r)=\sum_{j=1}^L D_j(r)n^{-j}
          +O_L\left(\frac{r^{L+2}}{n^{L+1}}\right)
       =\sum_{j=1}^L t^{2j}D_j(r)+O_L(t^L).
\end{equation}
To verify this, rewrite the defect exactly as
\[
 \Delta_n(r)=\frac{n-r}{2}\log(1+r/n)-\frac r2
            +\sum_{h=0}^{r-1}\log(1+(r-2h)/n).
\]
Expand the first logarithm through degree $L+1$ and each logarithm in
the sum through degree $L$. Their arguments are within $O(t)$ of one.
The usual logarithmic remainders, together with the $n,r$ prefactors
and the $r$ terms in the sum, give $O(r^{L+2}/n^{L+1})$.
The statement at $r=0$ is exact. Collecting coefficients gives
\eqref{spb:ep:eq:Dcrit}.

\subsection{Centered moments and a finite weighted generator}
Put
\begin{equation}\label{spb:ep:eq:Ucenter}
 U=t(X-\lambda)=tX-c.
\end{equation}
Let $B_{d,k}^{\ge2}$ count set partitions of $d$ labeled elements into
$k$ blocks, with no singleton blocks. They satisfy
\begin{equation}\label{spb:ep:eq:assocStirling}
 B_{0,0}^{\ge2}=1,\qquad
 B_{d,k}^{\ge2}=kB_{d-1,k}^{\ge2}+(d-1)B_{d-2,k-1}^{\ge2},
\end{equation}
with out-of-range indices zero. The recurrence separates the block
containing the last label according as it has more than two elements or
has two elements.
The centered Poisson generating function is
\[
 \E e^{z(X-\lambda)}=\exp(\lambda(e^z-1-z)).
\]
Coefficient comparison gives the exact polynomial identity
\begin{equation}\label{spb:ep:eq:centralmoment}
 \E U^d=\sum_{k=0}^{\lfloor d/2\rfloor}
           B_{d,k}^{\ge2}c^k t^{d-k}.
\end{equation}
Thus only integer powers of $t$ occur, and their least possible exponent
is $\lceil d/2\rceil$. In particular,
$\E|U|^{2h}=O_h(t^h)$ for each fixed $h$.

For the chosen expansion order $K$, define
\begin{equation}\label{spb:ep:eq:VK}
 V_K(t,u)=\frac{3c^2}{4}
            +\sum_{j=1}^{K+1}t^{2j}D_j((c+u)/t).
\end{equation}
It is a polynomial with no negative powers of $t$ because
$\deg D_j\le j+1$. Its constant term vanishes and
\begin{equation}\label{spb:ep:eq:VK0}
 V_K(0,u)=-\frac{3cu}{2}-\frac{3u^2}{4}.
\end{equation}
Assign weight two to $t$ and weight one to $u$.
Let $P_K(t,u)$ be $e^{V_K(t,u)}$ truncated to weighted degree at most
$2K+1$. Every monomial of $V_K$ has positive weight, so it suffices to
truncate the finite sum
\begin{equation}\label{spb:ep:eq:PK}
 \sum_{m=0}^{2K+1}\frac{V_K(t,u)^m}{m!}.
\end{equation}
Replace every $u^d$ in $P_K$ by the right side of
\eqref{spb:ep:eq:centralmoment}, and retain powers through $t^K$.
This finite procedure defines $C_{1,0},\ldots,C_{1,K}$.
It is different from the small-deformation generator of the previous section.

\subsection{Remainder and parity proof}
\begin{proof}[Proof of Theorem~\ref{spb:ep:thm:critical}]
On $\mathcal C_n$, one has
$|U|\le t\lambda^{3/4}=O(t^{1/4})\to0$.
Take $L=K+1$ in \eqref{spb:ep:eq:criticalDrem}. The exponent error is
$O(t^{K+1})$, and $V_K(t,U)$ is uniformly bounded on this event.
Therefore
\begin{equation}\label{spb:ep:eq:criticalexprem}
 e^{\Delta_n(X)+3c^2/4}
       =e^{V_K(t,U)}+O_K(t^{K+1})
       \quad\hbox{on }\mathcal C_n.
\end{equation}

Set $s=\sqrt t$ solely to apply ordinary multivariate Taylor's theorem to
$e^{V_K(s^2,u)}$ at $(0,0)$. Its total-degree $2K+1$ Taylor polynomial is
$P_K(s^2,u)$. The remainder is
$O_K((|s|+|u|)^{2K+2})$ in a fixed neighborhood.
Averaging this bound, using \eqref{spb:ep:eq:centralmoment} and
$(a+b)^m\le2^{m-1}(a^m+b^m)$, gives $O_K(t^{K+1})$.
The parity-restricted expectation costs at most a factor two.
The extra odd weighted degree is important: a Taylor argument stopped at
$2K$ would give only an absolute $O(t^{K+1/2})$ bound by this method.

The polynomial tail estimate and \eqref{spb:ep:eq:criticaltail} allow us to replace
the central expectation of $P_K$ by its unrestricted Poisson expectation.
For every fixed polynomial $Q(X)$ the exact parity filter says
\[
 2\E[Q(X)\1_{\{X\equiv n\ (2)\}}]
    =\E Q(X)+(-1)^n\E[(-1)^XQ(X)].
\]
For $Q(X)=X^j$, the last expectation is
$e^{-2\lambda}T_j(-\lambda)$.
A monomial $t^iU^d=t^i(tX-c)^d$ is a finite combination of
$t^{i+j}X^j$. Since $t^jT_j(-c/t)$ remains bounded, the parity contribution
from $P_K$ is $O_K(e^{-2\lambda})$.
Finally \eqref{spb:ep:eq:centralmoment} produces only integer powers of $t$;
those beyond $K$ are $O_K(t^{K+1})$.
Combining these statements with \eqref{spb:ep:eq:criticalexact} proves the
expansion and the claimed absence of odd quarter powers.
The displayed coefficients result from the finite generator.
\end{proof}

\subsection{The first coefficient without symbolic expansion}
There is a short independent way to see the first correction.
Put $Z=tX=c+U$ and $f(z)=e^{-3z^2/4}$.
On the bulk, the first two defect polynomials give
\[
 e^{\Delta_n(X)}=f(Z)\{1+t(Z+Z^3/4)+O(t^2)\}.
\]
Since $\E U=0$ and $\E U^2=ct$, the centered Taylor calculation gives
\[
 \E f(Z)=f(c)+\frac{ct}{2}f''(c)+O(t^2).
\]
Here a fourth-order Taylor remainder and the exact third central moment
justify the $O(t^2)$ error; an absolute third-moment bound alone would be
insufficient. The same argument applied to $f(Z)(Z+Z^3/4)$ gives its
expectation equal to its value at $c$ plus $O(t)$.
The tails have already been controlled globally.
Dividing by $f(c)$, the coefficient is
\[
 c+\frac{c^3}{4}+\frac c2\left(-\frac32+\frac{9c^2}{4}\right)
       =\frac c4+\frac{11c^3}{8}.
\]
This identifies the missing cancellation without a symbolic black box.
It also explains why adding the observed extra $1/(8c)$ changes the true
refinement while leaving the leading equivalent untouched.

\section{The inactive vertex distribution}
\label{sec:marked}
This entire section assumes fixed $p\ge2$.
Let $R_n$ be the inactive-vertex count from Section~\ref{sec:exact}, let
$a\equiv n\pmod q$, and let $Q_{\nu,a}$ be the distribution of
$\Pois(\nu)$ conditioned on the residue class $a$ modulo $q$.
For probability measures on this lattice we use
\[
 \TV(P,Q)=\frac12\sum_r|P(r)-Q(r)|.
\]
The exact Poisson identity shows that the mass of $R_n$ is proportional to
\begin{equation}\label{eq:markedweight}
 \frac{\lambda^r}{r!}e^{\Delta_n(r)},
 \qquad 0\le r\le n,\quad r\equiv a\pmod q.
\end{equation}
All constants $A,b,\gamma_1,\gamma_2,\gamma_3$ retain the meanings in
\eqref{eq:firstuniv} and \eqref{eq:D2}.

\subsection{The two Poisson references}
\begin{theorem}\label{thm:markedmain}
There is a unique positive solution $\mu$ to
\begin{equation}\label{eq:mu}
 \log(\mu/\lambda)=-\frac{2A\mu+b}{n}.
\end{equation}
In terms of the principal positive branch of the Lambert function,
\begin{equation}\label{eq:muW}
 \mu=\frac{n}{2A}W\left(\frac{2A\lambda}{n}e^{-b/n}\right)
 =\lambda-\frac{2A\lambda^2+b\lambda}{n}
                   +O_p(\lambda^3/n^2).
\end{equation}
The mean and variance satisfy
\begin{align}
 \E R_n&=\lambda-\frac{2A\lambda^2+b\lambda}{n}
                      +O_p(\lambda^3/n^2),\label{eq:mean}\\
 \Var R_n&=\lambda-\frac{4A\lambda^2+b\lambda}{n}
                      +O_p(\lambda^3/n^2).\label{eq:variance}
\end{align}
Moreover,
\begin{align}
 \TV(\mathcal L(R_n),Q_{\lambda,a})
       &\sim A\sqrt{2/\pi}\,\frac{\lambda^{3/2}}{n},\label{eq:TVlambda}\\
 \TV(\mathcal L(R_n),Q_{\mu,a})
       &\sim A\sqrt{2/(\pi e)}\,\frac{\mu}{n}.\label{eq:TVmu}
\end{align}
The estimates and equivalents hold uniformly over the residue classes.
\end{theorem}

Uniqueness in \eqref{eq:mu} follows because
$\log\mu+2A\mu/n$ is strictly increasing from $-\infty$ to $+\infty$.
Equation \eqref{eq:muW} follows from $W(z)e^{W(z)}=z$;
its expansion uses $\lambda/n\to0$.
Since $\lambda^2/n\to0$, the remainder in \eqref{eq:mean} and
\eqref{eq:variance} is smaller even than the displayed $\lambda/n$ terms.
In particular,
\begin{equation}\label{eq:fano}
 \frac{\Var R_n}{\E R_n}
       =1-\frac{2A\lambda}{n}+O_p(\lambda^2/n^2).
\end{equation}

The reference in both total-variation statements is conditioned.
An unconditioned Poisson distribution gives mass approaching $1-1/q$ to
the complement of the support of $R_n$, so its total variation from
$R_n$ is at least $1-1/q+o(1)$. A Gaussian limit after scaling and a
small total-variation distance on the original integer lattice are
therefore different assertions.

\subsection{A density expansion with weighted remainder}
Define the monic second Poisson--Charlier polynomial by
\begin{equation}\label{eq:charlier}
 \mathcal C_2(r;\mu)=(r-\mu)^2-r.
\end{equation}
This normalization is $\mu^2$ times the degree-two Charlier polynomial
in the generating-function convention of DLMF 18.23.5~\cite{dlmf}.
Let
\begin{equation}\label{eq:K2}
 K_2(r;\mu)=D_2(r)+\frac{A^2}{2}\mathcal C_2(r;\mu)^2.
\end{equation}
Write $\E_\mu$ for ordinary Poisson expectation and $\E_Q$ for expectation
under $Q_{\mu,a}$.

\begin{theorem}\label{thm:weighted}
On the support of $Q_{\mu,a}$ define
$g_n(r)=\Pp(R_n=r)/Q_{\mu,a}(r)$, taking $g_n(r)=0$ for $r>n$.
Then
\begin{align}\label{eq:secondensity}
 g_n(r)={}&1-\frac{A}{n}\mathcal C_2(r;\mu)\nonumber\\
 &+\frac{K_2(r;\mu)-\E_\mu K_2(X;\mu)}{n^2}
   +\epsilon_n(r),
\end{align}
where, for $s=0,1,2$,
\begin{equation}\label{eq:weightedrem}
 \E_Q[|X-\mu|^s|\epsilon_n(X)|]
                       =O_p(\mu^{4+s/2}/n^3).
\end{equation}
Consequently,
\begin{equation}\label{eq:firstdensity}
 \E_Q\left|g_n(X)-1+\frac A n\mathcal C_2(X;\mu)\right|
                          =O_p(\mu^{5/2}/n^2).
\end{equation}
For each fixed $L>0$, uniformly at allowed points with
$|r-\mu|\le L\sqrt\mu$,
\begin{equation}\label{eq:localdensity}
 g_n(r)=1-\frac A n\mathcal C_2(r;\mu)
                         +O_{p,L}(\mu^{5/2}/n^2).
\end{equation}
\end{theorem}

\begin{proof}
Write $\beta_0=\log(\mu/\lambda)=-(2A\mu+b)/n$.
Changing the reference Poisson parameter subtracts $\beta_0r$ from the
exponent in \eqref{eq:markedweight}, up to a normalizing constant.
The identity $q/2+b=A$ gives the exact cancellation
\begin{equation}\label{eq:cancel}
 \frac{D_1(r)}n-\beta_0r
       =\frac{A\mu^2}{n}-\frac A n\mathcal C_2(r;\mu).
\end{equation}
The constant disappears after normalization. On $r\le2\mu$, logarithmic
Taylor expansion gives
\[
 \Delta_n(r)=\frac{D_1(r)}n+\frac{D_2(r)}{n^2}
                                  +O_p(r^4/n^3).
\]
Since $\mu/\lambda\to1$ and $\mu^2/n\to0$, the residual exponent is
uniformly $o(1)$ on this interval:
\begin{equation}\label{eq:residualV}
 V(r)=-\frac A n\mathcal C_2(r;\mu)
              +\frac{D_2(r)}{n^2}+O_p(r^4/n^3).
\end{equation}
Its exponential is
$1-A\mathcal C_2/n+K_2/n^2$ plus a remainder.

We give the weighted estimates explicitly. For $Y=X-\mu$, the centered
Poisson cumulant generating function shows that
$\E_\mu Y^{2h}=O_h(\mu^h)$ for $\mu\ge1$.
Cauchy--Schwarz and these even moments imply, for $s=0,1,2$,
\begin{align*}
 \E_\mu[|Y|^s|\mathcal C_2|^3]&=O_p(\mu^{3+s/2}),\\
 \E_\mu[|Y|^s|\mathcal C_2D_2|]&=O_p(\mu^{4+s/2}),\\
 \E_\mu[|Y|^sD_2^2]&=O_p(\mu^{6+s/2}),\\
 \E_\mu[|Y|^sX^4]&=O_p(\mu^{4+s/2}).
\end{align*}
For example $\mathcal C_2=Y^2-Y-\mu$ and $X=\mu+Y$ reduce each assertion
to a fixed collection of central moments. Conditioning changes upper
bounds only by a fixed factor because
$\Pp_\mu(X\equiv a\pmod q)\to1/q$.
The four displayed bounds control respectively the cubic exponential
term, the cross term, the square of $D_2$, and the exponent remainder.
The $D_2^2$ contribution carries $n^{-4}$ and is absorbed into
$O(\mu^{4+s/2}/n^3)$ because $\mu^2/n\to0$.

Tails beyond $2\mu$ are exponentially small after any fixed polynomial
weight. For the exact law, this follows from the original
$e^{\Delta_n(r)}\le1$ domination under $\Pois(\lambda)$, the normalizer
tending to one, and $2\mu/\lambda\to2$.
For the reference it is an ordinary Chernoff bound.
Thus no global pointwise estimate on the changed-reference density is
being assumed.

It remains to check normalization. The exact identities
$\E_\mu\mathcal C_2=0$ and
$\E_\mu\mathcal C_2^2=2\mu^2$ give
$\E_\mu K_2=O_p(\mu^3)$.
The fixed-polynomial residue correction is exponentially small by
\eqref{eq:filter}. Dividing by the normalizer subtracts
$\E_\mu K_2/n^2$ at second order. Its cross term with
$\mathcal C_2/n$ has weighted size
$O(\mu^{4+s/2}/n^3)$, and the normalizer-square terms have size
$O(\mu^{6+s/2}/n^4)$, again absorbed by the same bound.
This proves \eqref{eq:secondensity}--\eqref{eq:weightedrem}.

Finally, $D_2-\E_\mu D_2$ has $L^1$ norm $O_p(\mu^{5/2})$,
while $\mathcal C_2^2-\E_\mu\mathcal C_2^2$ has $L^1$ norm
$O(\mu^2)$. These follow by expanding in $Y$ and using its central moments.
They prove \eqref{eq:firstdensity}; its third-order remainder is smaller
because $\mu^{3/2}/n\to0$.
On $|r-\mu|\le L\sqrt\mu$, the corresponding pointwise centered bounds
hold, and the normalized third-order remainder is
$O_{p,L}(\mu^4/n^3)=o(\mu^{5/2}/n^2)$.
This proves \eqref{eq:localdensity}.
\end{proof}

\subsection{Moments and their second correction}
The polynomial identities needed to extract the moments are
\begin{align}
 \E_\mu\mathcal C_2&=0,&
 \E_\mu[(X-\mu)\mathcal C_2]&=0,\label{eq:charlierorth}\\
 \E_\mu[((X-\mu)^2-\mu)\mathcal C_2]&=2\mu^2,&
 \E_\mu[(X-\mu)\mathcal C_2^2]&=4\mu^2,\nonumber\\
 \E_\mu[((X-\mu)^2-\mu)\mathcal C_2^2]
        &=8\mu^3+8\mu^2.&&\nonumber
\end{align}
They can be verified directly by the factorial moment rule.
More generally, Poisson summation by parts gives for any polynomial $f$
\begin{align}
 \Cov_\mu(X,f(X))&=\mu\E_\mu[f(X+1)-f(X)],\label{eq:parts1}\\
 \Cov_\mu((X-\mu)^2,f(X))
  &=\mu\E_\mu[f(X+1)-f(X)]\nonumber\\
  &\quad+\mu^2\E_\mu[f(X+2)-2f(X+1)+f(X)].\label{eq:parts2}
\end{align}
For example, $\E[Xf(X)]=\mu\E f(X+1)$ gives the first identity;
using $\E[(X)_2f(X)]=\mu^2\E f(X+2)$ gives the second.

Multiply \eqref{eq:secondensity} by $X-\mu$ or $(X-\mu)^2$.
The second-order polynomial terms are $O_p(\mu^3/n^2)$ by
\eqref{eq:charlierorth}--\eqref{eq:parts2}.
The weighted errors are $O_p(\mu^{9/2}/n^3)$ and
$O_p(\mu^5/n^3)$, respectively, both $o(\mu^3/n^2)$.
The small centered mean squared is absorbed in the variance error.
Consequently
\begin{equation}\label{eq:meanmu}
 \E R_n=\mu+O_p(\mu^3/n^2),\qquad
 \Var R_n=\mu-\frac{2A\mu^2}{n}+O_p(\mu^3/n^2).
\end{equation}
Substitution of \eqref{eq:muW} proves \eqref{eq:mean} and
\eqref{eq:variance}.

A useful second-order form is obtained by setting $\theta=z\,d/dz$ and
\begin{align}\label{eq:L2}
 L_2(z)={}&(\gamma_3+2A^2)z^3
 +(3\gamma_3-\gamma_2+A^2+2Ab)z^2\nonumber\\
 &+(\gamma_3-\gamma_2-\gamma_1+b^2/2)z.
\end{align}
Then the same weighted calculation, retaining its polynomial terms, gives
\begin{align}
 \E R_n={}&\lambda-\frac{2A\lambda^2+b\lambda}{n}
            +\frac{(\theta L_2)(\lambda)}{n^2}
            +O_p(\lambda^{9/2}/n^3),\label{eq:mean2}\\
 \Var R_n={}&\lambda-\frac{4A\lambda^2+b\lambda}{n}
            +\frac{(\theta^2L_2)(\lambda)}{n^2}
            +O_p(\lambda^5/n^3).\label{eq:var2}
\end{align}
These are conservative weighted remainder bounds.
An independent way to verify the coefficient is to write
$D_1=-A(X)_2-bX$ and use
\[
 L_2=\E D_2+\tfrac12\Var(D_1),\quad
 \Var((X)_2)=4\lambda^3+2\lambda^2,\quad
 \Cov((X)_2,X)=2\lambda^2.
\]
It yields exactly \eqref{eq:L2}. In the two highlighted cases,
\begin{align*}
 p=2:\quad &L_2(z)=12z^3+\frac{44}{3}z^2+\frac43z,\\
 p=3:\quad &L_2(z)=\frac{4967}{96}z^3
                      +\frac{4515}{64}z^2+\frac{1161}{128}z.
\end{align*}

\subsection{Sharp total variation constants}
\begin{proof}[Proof of \eqref{eq:TVlambda} and \eqref{eq:TVmu}]
Expansion at the original reference and division by its normalizer give
\begin{equation}\label{eq:originaldensity}
 \frac{\Pp(R_n=r)}{Q_{\lambda,a}(r)}
 =1+\frac{D_1(r)-\E_{Q_{\lambda,a}}D_1}{n}+\rho_n(r),
 \qquad \E_{Q_{\lambda,a}}|\rho_n|=O_p(\lambda^4/n^2).
\end{equation}
The same local expansion and global tail proof as in Section~\ref{sec:subcritical}
justifies this $L^1$ statement, including the cutoff.
Under $Q_{\lambda,a}$, write $Z_\lambda=(X-\lambda)/\sqrt\lambda$.
Then $Z_\lambda$ converges to a standard normal $Z$, uniformly over residues,
and its fixed even moments are bounded.
To see convergence, apply the residue filter to the characteristic function
at $t/\sqrt\lambda$. The zero root gives the ordinary Poisson central
limit, and every nonzero root decays exponentially for bounded $t$.
Moment bounds follow from the fixed-polynomial filter.

Expansion of $D_1$ about $\lambda$ now gives
\[
 \frac{D_1(X)-\E_{Q_{\lambda,a}}D_1}{\lambda^{3/2}}
                    \longrightarrow -2AZ
\]
in distribution and in absolute first moment, by uniform integrability.
Since $\lambda^{5/2}/n\to0$ for $p\ge2$, the error in
\eqref{eq:originaldensity} is negligible relative to
$\lambda^{3/2}/n$. Thus
\[
 \TV(\mathcal L(R_n),Q_{\lambda,a})
     \sim\frac{2A\lambda^{3/2}}{2n}\E|Z|
     =A\sqrt{2/\pi}\,\frac{\lambda^{3/2}}n.
\]

For the tilted reference, \eqref{eq:firstdensity} gives
\[
 \TV(\mathcal L(R_n),Q_{\mu,a})
 =\frac{A}{2n}\E_{Q_{\mu,a}}|\mathcal C_2(X;\mu)|
                                    +O_p(\mu^{5/2}/n^2).
\]
The same conditional central limit and uniform integrability imply
$\mathcal C_2(X;\mu)/\mu\to Z^2-1$ in distribution and absolute first
moment. If $\phi$ is the standard normal density, integration by parts
on $[0,1]$ gives $\E|Z^2-1|=4\phi(1)$.
Since $\mu^{3/2}/n\to0$, the error is negligible, and the leading constant is
$2A\phi(1)=A\sqrt{2/(\pi e)}$.
\end{proof}

\subsection{All fixed marked orders and graph consequences}
The finite polynomials $U_j$ in \eqref{eq:U} give an algebraic logarithm
by the recurrence
\begin{equation}\label{eq:Lrec}
 L_j(z)=U_j(z)-\frac1j\sum_{k=1}^{j-1}kL_k(z)U_{j-k}(z).
\end{equation}
Equivalently, as a formal series in $w$,
\[
 \log\left(1+\sum_{j\ge1}U_j(z)w^j\right)
                      =\sum_{j\ge1}L_j(z)w^j.
\]
The formal coefficient of the $h$th cumulant of $R_n$ is then
\begin{equation}\label{eq:cumulants}
 \lambda+\sum_{j\ge1}\frac{(\theta^hL_j)(\lambda)}{n^j}.
\end{equation}
This is an effective all-fixed-order generator, with the following precise
remainder interpretation. Multiply the finite density expansion through
$J$ by each required fixed power $r^d$, apply the polynomial filter, and
form the finite quotients and cumulant combinations. The unnormalized
weighted error is bounded by
\[
 O_{p,J,d}(\lambda^{2J+2+d}/n^{J+1}).
\]
For a desired fixed algebraic order, take $J$ sufficiently large and then
collect powers after setting $\lambda=c_p n^{1/q}$.
This procedure justifies \eqref{eq:cumulants} to that order without
assuming any general degree cancellation in $L_j$, and without
differentiating an uncontrolled asymptotic remainder.

For clarity, an exact marked generating function can be written before
any approximation:
\[
 \E e^{sR_n}=
 \frac{\sum_{r\equiv a\ (q),\,r\le n}(\lambda e^s)^r e^{\Delta_n(r)}/r!}
      {\sum_{r\equiv a\ (q),\,r\le n}\lambda^r e^{\Delta_n(r)}/r!}.
\]
Its finite derivatives explain the operator $\theta$ in the coefficient
rule. The error control is supplied by the weighted argument above.

Finally, if $K_n$ and $V_n$ are the numbers of active and total vertices,
\eqref{eq:rdef} gives the exact relations
\[
 \E K_n=\frac{n-\E R_n}{q},\quad
 \Var K_n=\frac{\Var R_n}{q^2},\quad
 \E V_n=\frac{n+p\E R_n}{q},\quad
 \Var V_n=\frac{p^2\Var R_n}{q^2}.
\]
The two vertex counts are affinely determined by the same defect.
Their centered Gaussian limits follow from the conditional Poisson limit
and either total-variation approximation. None of the statements in
this section is asserted for $p=1$.

\section{Finite smooth models and controlled inverses}
\label{sec:inverse}
This section covers every fixed integer $p\ge1$, using the counting theorem
appropriate to that $p$. Its subject is inversion of explicit finite smooth
models. We do not differentiate the discrete sequence, nor an unspecified
error in its asymptotic expansion.

\subsection{Unified notation and model derivatives}
Put
\begin{equation}\label{eq:inverseparams}
 \beta=1/q,\qquad \kappa=p/q,\qquad \alpha=c_p,\qquad
 d_p=\begin{cases}3e/8,&p=1,\\0,&p\ge2,\end{cases}
\end{equation}
and use the smooth positive function
\begin{equation}\label{eq:smoothM}
 M_p(x)=q^{-1}(x/q)^{\kappa x}\exp(\alpha x^\beta-d_p),
 \qquad x>0.
\end{equation}
For a fixed $K\ge0$, define the finite sum
\begin{equation}\label{eq:QK}
 Q_K(x)=\sum_{j=0}^K C_{p,j}(\alpha)x^{-j/q},\qquad
 \nu_K=(K+1)/q.
\end{equation}
For $p=1$, these are the separate critical coefficients from
Theorem~\ref{thm:critical}. Both counting theorems give
\begin{equation}\label{eq:modelvalue}
 a_p(n)/M_p(n)=Q_K(n)+O_{p,K}(n^{-\nu_K}).
\end{equation}
Let
\begin{equation}\label{eq:Fmodel}
 F_K(x)=\kappa x\log(x/q)+\alpha x^\beta-\log q-d_p+\log Q_K(x)
\end{equation}
where $x$ is sufficiently large that $Q_K(x)>0$.

\begin{lemma}\label{lem:model}
For each fixed $p,K$, there exists $X_{p,K}>e$ such that for
$x\ge X_{p,K}$,
\begin{equation}\label{eq:modelderivative}
 Q_K(x)>1/2,\qquad F_K'(x)\ge c_p'\log x>0,\qquad
             0<F_K''(x)\le C_{p,K}/x.
\end{equation}
The function $F_K$ has a unique smooth inverse $u_K$ on its eventual range,
and at integers
\begin{equation}\label{eq:logcomparison}
 \log a_p(n)=F_K(n)+O_{p,K}(n^{-\nu_K}).
\end{equation}
For fixed $L\ge K$ and each fixed derivative order $h\ge0$,
\begin{equation}\label{eq:coherentF}
 (F_L-F_K)^{(h)}(x)=O_{p,K,L,h}(x^{-\nu_K-h}).
\end{equation}
\end{lemma}
\begin{proof}
Let $\delta=1/2$ for $p=1$ and $\delta=(p-1)/q$ for $p\ge2$.
The vanishing coefficients already proved imply
$(Q_K-1)^{(h)}(x)=O_{p,K,h}(x^{-\delta-h})$;
if $Q_K=1$ the left side is zero. This is ordinary differentiation of a
known finite expression. In particular $Q_K>1/2$ eventually, and
$(\log Q_K)^{(h)}=O(x^{-\delta-h})$.
Thus
\begin{align}
 F_K'(x)&=\kappa(\log(x/q)+1)+\alpha\beta x^{\beta-1}
                                      +O(x^{-\delta-1}),\label{eq:Fprime}\\
 F_K''(x)&=\kappa/x+\alpha\beta(\beta-1)x^{\beta-2}
                                      +O(x^{-\delta-2}).\label{eq:Fsecondmodel}
\end{align}
The positive $\kappa/x$ term dominates the second derivative, proving
\eqref{eq:modelderivative}. Higher derivatives of order $h\ge2$ are
$O_{p,K,h}(x^{1-h})$. Taking a logarithm in \eqref{eq:modelvalue}, using
$Q_K>1/2$, proves \eqref{eq:logcomparison}.
Finally $Q_L-Q_K$ is a finite sum starting at $x^{-\nu_K}$, so its
fixed derivatives and those of the logarithmic difference have the stated
bounds. This proves \eqref{eq:coherentF}.
\end{proof}

\subsection{A Lambert core and finitely many Newton steps}
Write $Y=\log y$ and set
\begin{equation}\label{eq:Lambertcore}
 Z=Y+\log q+d_p,\qquad
 N=N(Y)=\frac{qZ}{pW(Z/p)},
\end{equation}
where $W$ is the positive real branch.
This solves exactly
\begin{equation}\label{eq:coreeq}
 \kappa N\log(N/q)=Z.
\end{equation}
Indeed, $N/q=\exp(W(Z/p))$ and $W(Z/p)e^{W(Z/p)}=Z/p$.
For fixed $K,J$, begin at $x_0=N(Y)$ and make $J$ Newton steps
\begin{equation}\label{eq:newton}
 x_{j+1}=x_j-\frac{F_K(x_j)-Y}{F_K'(x_j)},\qquad 0\le j<J.
\end{equation}
This is an explicit finite composition of $W$, real powers, logarithms,
exponentials, and arithmetic. There is no hidden root solve in $x_J$.

\begin{theorem}\label{thm:newton}
For every fixed $p,K,J$, all iterates in \eqref{eq:newton} are well-defined
for sufficiently large $Y$, and
\begin{equation}\label{eq:monotoneNewton}
 u_K(Y)\le x_{j+1}\le x_j\le N(Y).
\end{equation}
Their error obeys
\begin{equation}\label{eq:newtonerror}
 0\le x_J-u_K(Y)
   \le C_{p,K,J}\,
       \frac{N^{1-(1-\beta)2^J}}{(\log N)^{2^{J+1}-1}}.
\end{equation}
\end{theorem}
\begin{proof}
At the starting point,
$F_K(N)-Y=\alpha N^\beta+\log Q_K(N)>0$ eventually.
On the other hand $F_K(N/2)<Y$ for large $Y$, because the change in the
core term is of order $N\log N$, much larger than $N^\beta$.
Thus $u=u_K(Y)\in(N/2,N)$.
By the mean-value theorem and \eqref{eq:modelderivative},
\begin{equation}\label{eq:e0}
 e_0=N-u=O_{p,K}(N^\beta/\log N).
\end{equation}
On $[u,N]$, $F_K$ is increasing and convex. The tangent line at $x\ge u$
meets the line of height $Y$ between $u$ and $x$; this proves the
monotonicity and all domain assertions in \eqref{eq:monotoneNewton}.
Taylor's theorem, expanded at $x_j$, gives
\[
 e_{j+1}=\frac{F_K''(\xi_j)}{2F_K'(x_j)}e_j^2
                \le C\frac{e_j^2}{N\log N}.
\]
The power of $N$ therefore changes by $a\mapsto2a-1$ and the power in
the denominator by $b\mapsto2b+1$. Starting from $(a,b)=(\beta,1)$ in
\eqref{eq:e0} yields \eqref{eq:newtonerror}.
\end{proof}

The first Newton error is $O(N^{2\beta-1}/\log^3N)$.
At criticality $\beta=1/2$, so the first two Newton errors are respectively
\begin{equation}\label{eq:criticalNewton}
 O(\log^{-3}N),\qquad O(N^{-1}\log^{-7}N).
\end{equation}
The constant offset in the critical core is essential:
$N=2(Y+\log2+3e/8)/W(Y+\log2+3e/8)$.

\subsection{Arbitrary fixed index accuracy}
\begin{corollary}\label{cor:index}
Given any fixed $M>0$, choose integers $K,J\ge0$ with
\begin{equation}\label{eq:chooseKJ}
 (K+1)/q\ge M,\qquad (1-\beta)2^J\ge M+1.
\end{equation}
Let $h_{K,J}(Y)=x_J(Y)$. Then on the exact range of the count sequence,
\begin{equation}\label{eq:indexrecover}
 h_{K,J}(\log a_p(n))=n+O_{p,K,J}(n^{-M}/\log n).
\end{equation}
In particular, its nearest integer is $n$ for all sufficiently large $n$.
\end{corollary}
\begin{proof}
By \eqref{eq:logcomparison} and the derivative lower bound,
\[
 u_K(\log a_p(n))-n=O_{p,K}(n^{-\nu_K}/\log n).
\]
The leading growth also gives $N(\log a_p(n))\sim n$.
Combine this comparison with \eqref{eq:newtonerror} and
\eqref{eq:chooseKJ}. The logarithmic denominator in the Newton bound
is at least $\log N$, giving \eqref{eq:indexrecover}.
Eventually its absolute error is less than $1/2$.
\end{proof}

The inverse models are coherent. From \eqref{eq:coherentF} and the
mean-value theorem, for each fixed $L\ge K$,
\begin{equation}\label{eq:inversecoherent}
 u_L(Y)-u_K(Y)=O_{p,K,L}(N(Y)^{-\nu_K}/\log N(Y)).
\end{equation}
Increasing the model order therefore cannot change already resolved inverse
terms. All derivative bounds used here belong to explicit finite functions;
none is inferred by differentiating \eqref{eq:logcomparison}.

\subsection{A convergent local Lagrange generator}
For a fixed sufficiently large $N=N(Y)$, put
\begin{equation}\label{eq:HN}
 v=F_K(N)-Y,\qquad
 H_N(z)=\frac{z}{F_K(N+z)-F_K(N)},\qquad H_N(0)=\frac1{F_K'(N)}.
\end{equation}
The local Lagrange formula is
\begin{equation}\label{eq:Lagrange}
 u_K(Y)=N+\sum_{m\ge1}\frac{(-v)^m}{m!}
       \left.\frac{d^{m-1}}{dz^{m-1}}H_N(z)^m\right|_{z=0}.
\end{equation}
We verify a convergence disk and a quantitative truncation bound, rather
than using this as a merely formal identity.

\begin{proposition}\label{prop:Lagrange}
For fixed $p,K$ and all sufficiently large $Y$, series
\eqref{eq:Lagrange} converges. After retaining $m=1,\ldots,L$ its remainder is
\begin{equation}\label{eq:Lagrangerem}
 O_{p,K,L}\left(N\left(\frac{N^{\beta-1}}{\log N}\right)^{L+1}\right).
\end{equation}
\end{proposition}
\begin{proof}
Fix $0<\eta<1/4$. The explicit finite expression $F_K(N+z)$ is holomorphic
on a neighborhood of $|z|\le\eta N$ for large $N$: $N+z$ stays in the
right half-plane, and $Q_K(N+z)$ stays near one.
Uniformly on this disk,
\[
 F_K'(N+z)=\kappa\log N+O_{p,K,\eta}(1).
\]
Its real part is positive eventually. For two points in the disk, divide
their function-value difference by their point difference and integrate
the derivative along their connecting segment. The real part of the
result is positive, so the map is injective.

Write $G(z)=F_K(N+z)-F_K(N)$. Integrating the same derivative bound gives
\[
 |G(z)-\kappa(\log N)z|\le C|z|.
\]
Let $R_N=(\kappa\eta/4)N\log N$.
For large $N$, on $|z|=\eta N$ the perturbation is at most
$\tfrac14\kappa\eta N\log N$. For every $|w|<2R_N$,
\[
 |G(z)-\kappa(\log N)z-w|
      <\tfrac34\kappa\eta N\log N
      <|\kappa(\log N)z|.
\]
Rouch\'e's theorem gives exactly one solution of $G(z)=w$ inside this disk.
The derivative is nonzero there, so the inverse displacement is analytic
on the open image disk $|w|<2R_N$ and bounded by $\eta N$.
This strict radius margin permits Cauchy's estimate on $|w|=R_N$.

The Taylor coefficients of the inverse displacement are, by the local
Lagrange theorem~\cite{dlmf},
$\frac1{m!}(d^{m-1}/dz^{m-1})H_N(z)^m|_{z=0}$.
Since $v=O(N^\beta)$, one has $|v|/R_N\to0$.
Cauchy's bound and a geometric tail yield
\[
 O\left(N\frac{(|v|/R_N)^{L+1}}{1-|v|/R_N}\right),
\]
which proves \eqref{eq:Lagrangerem} and convergence at $w=-v$.
\end{proof}

The first two displacements in \eqref{eq:Lagrange} are
\begin{equation}\label{eq:firstLagrange}
 -\frac v{F_K'(N)}-\frac{F_K''(N)v^2}{2F_K'(N)^3}.
\end{equation}
The first leading displacement alone is
\begin{equation}\label{eq:firstdisplacement}
 -\frac{\alpha N^\beta}{\kappa(\log(N/q)+1)}.
\end{equation}
Discarding terms to get \eqref{eq:firstdisplacement} also discards their
accuracy. This shorter expression is not assigned the much smaller error
of a multistep Newton approximation.

\subsection{Exact integer thresholds and the rounding boundary}
\begin{theorem}\label{thm:threshold}
Fix $M>0$ and $K,J$ satisfying \eqref{eq:chooseKJ}.
For $h=h_{K,J}(\log y)$ there are constants $C,Y_0$, depending only on
$p,K,J$, such that for $y\ge e^{Y_0}$,
\begin{equation}\label{eq:ceiling}
 \left\lceil h-C\frac{h^{-M}}{\log h}\right\rceil
       \le T_p(y)\le
 \left\lceil h+C\frac{h^{-M}}{\log h}\right\rceil.
\end{equation}
The endpoint integers differ by at most one for all sufficiently large $y$.
If they coincide, that integer is the exact threshold. Otherwise the theorem
leaves two adjacent possibilities.
\end{theorem}
\begin{proof}
First use the exact model inverse $u=u_K(Y)$, $Y=\log y$.
For integers $m$ with $|m-u|\le2$, the error in
\eqref{eq:logcomparison} is bounded by $C_1u^{-\nu_K}$ and
$F_K'(x)\ge c_1\log u$ throughout this small neighborhood.
Choose
\[
 \varepsilon=\frac{2C_1u^{-\nu_K}}{c_1\log u}.
\]
For $m_+=\lceil u+\varepsilon\rceil$,
$F_K(m_+)-Y\ge2C_1u^{-\nu_K}$, hence
$\log a_p(m_+)\ge Y$.
For $m_-=\lceil u-\varepsilon\rceil-1$ one has strictly
$m_-<u-\varepsilon$, even if $u-\varepsilon$ is itself an integer.
It follows that $\log a_p(m_-)<Y$.
Exact monotonicity from Proposition~\ref{prop:monotone} yields
\[
 \lceil u-\varepsilon\rceil\le T_p(y)
                         \le\lceil u+\varepsilon\rceil.
\]
Replacing $u$ by $h$ uses \eqref{eq:newtonerror}; $u\sim h\sim N$ and
$\nu_K\ge M$ absorb both errors into the radius in \eqref{eq:ceiling}.
Twice that radius is eventually less than one, so its endpoint ceilings
can differ by at most one.
\end{proof}

\begin{remark}[Exact range versus arbitrary levels]
Nearest-integer recovery in Corollary~\ref{cor:index} assumes the input
is exactly $y=a_p(n)$. It does not give nearest-integer threshold recovery
for an arbitrary $y$. Even an extremely short real interval can straddle
an integer, in which case taking a ceiling cannot be settled by its width
alone. If actual numerical bounds and candidates are available, direct
integer evaluation can decide the remaining threshold comparison.
The present asymptotic theorem does not supply certified finite constants
or a certified onset for doing so.
\end{remark}

All onset thresholds and error constants in this inverse section are
existential. They could in principle be made effective from explicit
Taylor and concentration bounds, but the accompanying numerical inverse
experiments are diagnostics, not a certified finite-input algorithm.
The critical count theorem transfers to these inverses because its finite
models have the same derivative structure. It does not transfer the marked
probability theorem to $p=1$.

\section{Computation and source comparison}
\subsection{Exact arithmetic and independent checks}
The accompanying package separates three kinds of evidence:
\begin{enumerate}[leftmargin=*,itemsep=3pt]
 \item Integer and rational identities are checked with Python integers and
 \texttt{Fraction}. These include exact counts, defect polynomials,
 Poisson transforms, marked coefficients, and critical coefficients.
 \item Transcendental evaluations use optional \texttt{mpmath} arithmetic.
 They corroborate the theorems and expose finite-$n$ limitations; they
 are not certified interval enclosures.
 \item The proofs above provide asymptotic remainders and quantifiers.
 Neither fitted coefficients nor numerical agreement replaces those proofs.
\end{enumerate}
The exact core requires no symbolic algebra package. It computes finite
power sums by telescoping, the $H_j$ recurrence by rational polynomial
arithmetic, and Poisson transforms by the Touchard recurrence.
The critical generator separately performs weighted polynomial truncation
and exact centered-moment substitution.

Independent checks reconstruct polynomial expectations in the
falling-factorial basis. If a polynomial $H$ has degree at most $d$, then
\begin{equation}\label{spb:ep:eq:diffbasis}
 H(r)=\sum_{j=0}^d\Delta^jH(0)\binom rj,
 \qquad \E H(X)=\sum_{j=0}^d\Delta^jH(0)\frac{\lambda^j}{j!},
\end{equation}
where $\Delta H(r)=H(r+1)-H(r)$.
This is a distinct route from converting ordinary powers through Touchard
polynomials. Extra integer samples check the interpolated identities.
The public tests also multiply the formal generating function independently
at small sizes, check every residue for $p=1,2,3,4$ in numerical identity
tests, and test active input guards under optimized Python.

Program limits are deliberately finite resource safeguards, not restrictions
of the theorems. For example, the public generator caps the formal $H$ order
at 16, fractional order at 32 subject to that cap, and critical order at 8.
All guards use ordinary conditionals and exceptions, so Python's
\texttt{-O} flag cannot remove them. The precise limits and commands are
listed in the code README.

\subsection{Complete sums and a visible preasymptotic warning}
\label{spb:ep:sec:numerics}
For the complete-sum diagnostic, the original expression
\eqref{spb:ep:eq:counts} is evaluated as an exact integer. Its normalization uses
90-digit arithmetic. Independently, the transformed expression sums
\emph{every} allowed $r$ from zero through $n$, using logarithms and
log-gamma evaluations. No Poisson window and no asymptotic coefficients
enter either finite-sum comparison.

For $p\ge2$ let
\[
 A_K(n)=\sum_{k=0}^K C_{p,k}(c_p)n^{-k/q},\qquad
 S_J(n)=\sum_{j=0}^J U_j(\lambda)n^{-j}.
\]
These are different finite approximations. In particular, $S_{12}$ retains
many fractional powers beyond those in $A_{12}$.
Table~\ref{spb:ep:tab:counts} reports relative errors as approximation divided by
$\Rcal_p(n)$ minus one. The number of printed digits is deliberately smaller
than the working precision.

\begin{table}[htbp]
\centering\small
\begin{tabular}{@{}crrrr@{}}
\toprule
$p$ & $n$ & $\Rcal_p(n)$ & $A_{12}/\Rcal_p-1$ & $S_{12}/\Rcal_p-1$\\
\midrule
2 & 10000 & $0.5075173547743$ & $-9.64017\cdot10^{-9}$ & $+6.40647\cdot10^{-10}$\\
3 & 10000 & $0.1416076329487$ & $-0.6615364472$ & $+0.0012291163$\\
\bottomrule
\end{tabular}
\caption{Complete-sum diagnostics. Both parameterizations are asymptotically
valid at fixed order, but have sharply different errors here.}
\label{spb:ep:tab:counts}
\end{table}

The original and transformed evaluations agree to about 86 decimal places
in these cases. Such agreement checks the exact transformation numerically;
it does not certify the final transcendental digits by outward rounding.
Most importantly, the $p=3$ fractional truncation is still wrong by about
$66\%$ at $n=10000$. The shorter $t^2$ and $t^4$ approximations are worse.
Large constants such as $e^{9/2}$ delay the useful asymptotic regime.
The smaller error of $S_{12}$ at this input does not establish a general
optimal truncation principle.

For the critical case, the complete count at $n=10000$ gives
\[
 a_1(n)/M_1(n)\approx1.02446126439771172897.
\]
The half-power truncation through $C_{1,4}$ has relative error
approximately $-2.32029\cdot10^{-9}$.
The companion odd input $n=10001$ gives an error of the same size,
consistent with the common algebraic coefficients for both parities.
\begin{writenote}[The decimals recomputed, 7 October 2026]
The write recomputed Table~\ref{spb:ep:tab:counts} and the critical values
from the exact integers at 110 digits, with its own coefficients
$C_{p,k}$, $k\le12$, and transforms $U_j$, $j\le12$: every printed value is
the correct rounding of the recomputed one. Several are roundings rather than
truncations; the truncations at the printed length are $\Rcal_2=0.5075173547742\ldots$,
$A_{12}/\Rcal_2-1=-9.64016\ldots\cdot10^{-9}$,
$S_{12}/\Rcal_2-1=+6.40646\ldots\cdot10^{-10}$,
$\Rcal_3=0.1416076329486\ldots$, $A_{12}/\Rcal_3-1=-0.6615364471\ldots$,
$S_{12}/\Rcal_3-1=+0.0012291162\ldots$, $a_1(n)/M_1(n)=1.02446126439771172896\ldots$
and the relative error $-2.32028\ldots\cdot10^{-9}$ at $n=10000$
($-2.31971\ldots\cdot10^{-9}$ at $n=10001$). The integer and transformed
evaluations agree to $10^{-106}$ in the write's computation.
\end{writenote}

\subsection{Marked and inverse diagnostics}
At $n=10^9$ the reported TV distances divided by their proved leading
equivalents are approximately
\begin{center}
\begin{tabular}{@{}crr@{}}
\toprule
$p$ & Original conditioned reference & Tilted conditioned reference\\
\midrule
2 & $1.000111734$ & $1.000214479$\\
3 & $1.000578191$ & $0.999189636$\\
\bottomrule
\end{tabular}
\end{center}
These computations use central windows, with the window endpoints and
number of lattice points included in the machine-readable output.
Chernoff bounds control the difference between the window-normalized TV
and complete-law TV. Their high-precision evaluated bounds are tiny, but
are still not interval certificates. The moment outputs are explicitly
window diagnostics; a TV bound alone does not control unbounded moment
errors.

The Newton diagnostic has two modes. By default it sets $Y=F_K(n)$, so the
finite model root is known exactly to be $n$; this isolates iteration error.
An optional exact-target mode instead uses $Y=\log a_p(n)$ and therefore
includes model truncation error. The program makes exactly the requested
finite number of Newton steps and checks numerical domain conditions at
each step. Neither mode supplies a certified convexity onset or certifies
threshold rounding. These distinctions matter when interpreting a tiny
printed residual.
\begin{writenote}[The marked ratios recomputed, 7 October 2026]
With the exact law over a full central window (log-gamma values of the
summands) and the residue-conditioned Poisson references, the write obtains
the four ratios at $n=10^9$ as $1.00011173445\ldots$, $1.00021447898\ldots$
($p=2$) and $1.00057819116\ldots$, $0.999189635789\ldots$ ($p=3$): the printed
values are correct roundings (truncations $1.000214478$ and $0.999189635$ for
the two that differ).
\end{writenote}

\subsection{Observed primary formulas}
\label{spb:ep:sec:sources}
The exact finite sums and the three displayed OEIS expressions were
observed in primary-entry text on 5 October 2026.
The A360479 ordinary page was retrieved directly; A360592 and A360747
were obtained through indexed primary OEIS text when direct opening was
unreliable. Independent retrievals corroborated the mathematical formulas.
No verified sequence revision identifier or complete history export was
obtained. Global page footers and indexed modification dates do not
establish the current sequence revision.

A360592 identifies V\'aclav Kote\v{s}ovec as the sequence author, dated
13 February 2023. A360479 and A360747 identify Seiichi Manyama as sequence
author, dated 19 February 2023; their refinement formulas are attributed
to Kote\v{s}ovec on 19 and 20 February 2023 respectively.
These are observed source attributions, not a reconstructed publication
history. The inspected entry links were numerical tables or graphs rather
than proof articles.

Using the normalizations of this article, the observed critical multiplier
was
\[
 1+\left(\frac c4+\frac{11c^3}{8}+\frac1{8c}\right)t,
 \qquad c=\sqrt{e/2},\quad t=n^{-1/2}.
\]
For $p=2$, it was
\begin{align*}
 1&+\left(-\frac{13c^2}{6}+\frac1{8c}\right)t\\
  &+\left(\frac{169c^4}{72}-\frac{15c}{16}
                               +\frac{67}{384c^2}\right)t^2\\
  &+\left(-\frac{2197c^6}{1296}+\frac{7913c^3}{576}
                +\frac{3929}{2304}+\frac{497}{3072c^3}\right)t^3,
\end{align*}
with $c=e^{4/3}/3^{1/3}$ and $t=n^{-1/3}$.
For $p=3$, it was
\begin{align*}
 1&+\frac{t}{8c}
 +\left(-\frac{37c^2}{8}+\frac{67}{384c^2}\right)t^2\\
  &+\left(-\frac{205c}{64}+\frac{497}{3072c^3}\right)t^3\\
  &+\left(\frac{1369c^4}{128}+\frac{10721}{3072}
                      +\frac{218831}{1474560c^4}\right)t^4,
\end{align*}
with $c=e^{9/4}/\sqrt2$ and $t=n^{-1/4}$.
Each appeared as part of a formula using $\sim$.
Every one of these multipliers tends to one, so the formulas remain
correct if read solely as leading equivalents. Their intended refined
coefficients differ from \eqref{spb:ep:eq:critC1}, \eqref{spb:ep:eq:p2}, and
\eqref{spb:ep:eq:p3}. Appendix~\ref{spb:ep:app:discrepancies} lists the differences explicitly.
No explanation of how the discrepancies arose is assumed.
\begin{writenote}[7 October 2026]
The write compared these three multipliers with the formulas of the current
entries (revisions in Remark~\ref{spb:ep:rem:oeis}): after the normalizations
of this article they are exactly the displayed expressions (checked
symbolically, leading factors included).
\end{writenote}

\subsection{Bounded historical scope}
The source comparison establishes a specific coefficient audit, not a
universal priority result. Exact-number and formula-oriented searches did
not locate a proof of the observed refinements. Two inspected default-branch
catalogues in the public ProveIt repository and exact code searches did
not identify this endpoint family. Those are bounded negative checks;
they cannot rule out an unnamed general theorem or material elsewhere.

A potentially relevant 2013 note by Kote\v{s}ovec,
\emph{Interesting asymptotic formulas for binomial sums}, was bibliographically
identified at
\url{https://www.kotesovec.cz/math_articles/kotesovec_interesting_asymptotic_formulas.pdf},
but retrieval timed out and its theorem content was not compared.
Likewise the latest OEIS histories were not verified.
Consequently no assertion that all prior relevant work has been excluded
is made here. The correction follows from the exact finite sums and the
self-contained proofs, regardless of broader historical priority.

For classical background, the DLMF inverse-function and Charlier generating
function entries were checked directly. The Lambert-$W$ reference was
verified bibliographically and at its definition-level abstract; we do
not attribute an uninspected theorem in that paper to the present family.
The analytic inverse estimates needed here are proved in
Section~\ref{spb:ep:sec:inverse}.

\begin{remark}[{[write]} The OEIS entries, 7~October 2026]\label{spb:ep:rem:oeis}
The write read the three entries in the internal format on 7~October 2026.
\emph{A360592}, revision \#28 (17 February 2023), by V\'aclav Kote\v{s}ovec
(13 February 2023), ``G.f.: Sum\_\{k>=0\} (1 + k*x)\^{}k * x\^{}k.'', with the
finite sum, a b-file by Kote\v{s}ovec for $n=0,\ldots,760$, and the unsigned
formula line $a(n)\sim\exp(\exp(1/2)\sqrt{n/2}-3e/8)\,n^{n/2}/2^{n/2+1}\,
(1+((e^{1/2}+e^{-1/2})/2^{5/2}+11e^{3/2}/2^{9/2})/\sqrt n)$.
\emph{A360479}, revision \#40 (19 February 2023), by Seiichi Manyama
(19 February 2023), with a b-file by Kote\v{s}ovec ($n=0,\ldots,635$), a graph
of the asymptotic ratio, and the refinement through $n^{-1}$ signed by
Kote\v{s}ovec, 19 February 2023. \emph{A360747}, revision \#22 (20 February
2023), by Manyama (19 February 2023), with a b-file by Kote\v{s}ovec
($n=0,\ldots,600$), a graph, and the refinement through $n^{-1}$ signed by
Kote\v{s}ovec, 20 February 2023, ending ``see graph for more minor asymptotic
terms''. These are the latest revisions, all from February 2023, so the
formulas the source compared in Section~\ref{spb:ep:sec:sources} are those of
the current entries, and its attributions are as stated there.
\emph{Checks.} All terms of the three b-files equal the write's
evaluation of~\eqref{spb:ep:eq:counts}. Rewritten in this article's
normalizations, the three formula lines are exactly the multipliers displayed
in Section~\ref{spb:ep:sec:sources}, with leading factors $M_1(n)$, $M_2(n)$,
$M_3(n)$ (SymPy). The leading equivalents are therefore correct. The first
corrections are wrong: by Theorems~\ref{spb:ep:thm:critical}
and~\ref{spb:ep:thm:allp} the true first coefficients are $C_{1,1}$,
$C_{2,1}=-13c^2/6$ and $C_{3,1}=0$, and each entry adds $1/(8c_p)$. Against
the exact sums the write finds $(\text{entry multiplier}-a_p(n)/M_p(n))/t$
equal to $0.10722073\ldots$ at $n=10^{14}$ ($p=1$; $1/(8c)=0.10722048\ldots$),
$0.04752092\ldots$ at $n=10^{18}$ ($p=2$; $1/(8c_2)=0.04752160\ldots$) and
$0.01863213\ldots$ at $n=10^{24}$ ($p=3$; $1/(8c_3)=0.01863212\ldots$), the
full refined multipliers being used; so the entries' first corrections are
refuted by the proofs of this Part and, numerically, by the exact sums. No OEIS
edit is part of this work.
\end{remark}

\section{Further questions and boundaries of the results}
The finite-order method leaves several natural questions. Each requires
additional arguments beyond the theorems proved here.

\paragraph{Effective inverse constants.}
The local analytic bounds and Chernoff inequalities are explicit in form,
but this article has not optimized their constants or implemented a certified
finite onset. A useful next step is an interval-arithmetic inverse with a
provable error radius, followed by exact count comparisons for any
ambiguous adjacent thresholds.

\paragraph{A better moderate size approximation.}
The $p=3$ diagnostic suggests retaining more of the defect in exponential
form rather than immediately collecting fractional powers. The full
$H$ truncation is already much better at one displayed input, but this
alone proves no generally superior scheme. Uniform relative error bounds
for an appropriately resummed approximation would be more informative
than choosing a truncation by appearance.

\paragraph{Critical marked laws.}
At $p=1$ the leading quadratic tilt is order one in the partition function.
The count proof resums it, but the sharp marked-law argument in
Section~\ref{spb:ep:sec:marked} relies on $\lambda^2/n\to0$.
A separate centered analysis is needed to determine the best conditional
reference and sharp distances at criticality. Nothing in the inverse
transfer proves those probability results.
\begin{writenote}[Answered in Part~II, 7 October 2026]
Part~II (Report~235) carries out this centered analysis for the inactive
count: the sharp distances to both conditioned Poisson references
(Corollary~\ref{spb:ml:cor:Poisson}, with the constants of~\eqref{spb:ep:eq:TVlambda}
and~\eqref{spb:ep:eq:TVmu}), all fixed-order marked cumulants
(Theorem~\ref{spb:ml:thm:cumulants}), and the globally optimal
parity-conditioned binomial with its sharp constant
(Theorem~\ref{spb:ml:thm:global}). It does not determine a best reference
outside the conditioned binomial family.
\end{writenote}

\paragraph{Residue effects beyond algebraic order.}
The polynomial root filter displays exponentially small terms at nonzero
roots of unity. A remainder that is merely algebraically small can swamp
those terms. Establishing a secondary sector of the full deformed sum
would require a direct complex or oscillatory remainder analysis with
relative accuracy on that smaller scale. No such claim follows from
\eqref{spb:ep:eq:filter} alone.

\paragraph{Large order and growing parameters.}
All estimates hold at every chosen fixed order, with constants depending on
that order and on $p$. They give neither convergence of the infinite
asymptotic series nor optimal truncation. The scale $c_p$ grows rapidly
with $p$, so estimates uniform for a growing $p$ would need a fresh analysis.
There is also no continuous transition theorem joining $p=1$ to $p\ge2$;
the family considered here has integer $p$.

\paragraph{Sharper marked approximations.}
The tilted reference removes the leading location perturbation and leaves
a second Charlier correction. Higher signed Charlier corrections are
available algebraically. Turning them into positive probability
approximants with explicit distance bounds, or proving finer moderate
and large deviation estimates, would extend the present central laws.

\paragraph{Broader graph statistics.}
The inactive count determines the numbers of active and total vertices
through affine relations, but does not describe components, cycles,
indegrees, or repeated targets. Those observables require additional marks
in the directed-graph model. The exact size statistic and label convention
should be retained in any such extension.

The results established here are therefore concrete finite-order count,
probability, and inverse theorems. Their strengths are exact reduction,
global domination, explicit rational generators, and controlled error
transfer. Their scope does not depend on an unsupported claim about
exponentially small sectors or historical novelty.
\begin{writenote}[Added 7 October 2026, under Vladimir's standing rule of 4 October 2026]
Every open question of this Part is in this section, with what is missing;
the critical marked law is answered in Part~II within the conditioned binomial
family (note above). The write adds from the non-claims interval-certified values for the
tables and diagnostics of Section~\ref{spb:ep:sec:numerics}, which are
floating computations. No claim of the
source was found false; the OEIS refinements' first corrections are refuted
(Remark~\ref{spb:ep:rem:oeis}).
\end{writenote}

\appendix
\section{Coefficient and discrepancy reference}
\label{spb:ep:app:discrepancies}
For quick comparison, ``observed minus proved'' means the coefficient in the
normalized observed multiplier of Section~\ref{spb:ep:sec:sources} minus the
corresponding coefficient in this article. The parameter $c$ is always the
one belonging to the indicated $p$.
For $p=1$ the sole displayed difference is
\[
 [t]:\quad \frac1{8c}.
\]
For $p=2$ the differences through $t^3$ are
\begin{align*}
 [t]:&\quad \frac1{8c},\\
 [t^2]:&\quad \frac{67}{384c^2}-\frac{13c}{48},\\
 [t^3]:&\quad \frac{3929}{2304}+\frac{497}{3072c^3}
                                      +\frac{169c^3}{576}.
\end{align*}
For $p=3$ the differences through $t^4$ are
\begin{align*}
 [t]:&\quad \frac1{8c},\\
 [t^2]:&\quad \frac{67}{384c^2},\\
 [t^3]:&\quad \frac{497}{3072c^3}-\frac{37c}{64},\\
 [t^4]:&\quad \frac{10721}{3072}+\frac{218831}{1474560c^4}.
\end{align*}
These comparisons are about intended refinements, not the validity of
leading equivalence under the literal symbol $\sim$.

\section{Reproduction guide}
\label{spb:ep:app:reproduce}
The archive contains this PDF, editable TeX sources, the exact-arithmetic
code, optional high-precision diagnostics, checked outputs, a source guide,
and a SHA-256 manifest. No network access is needed for reproduction.

From an extracted archive, the standard-library exact checks are run with
\begin{verbatim}
python3 -B code/selftest.py
python3 -B -O code/selftest.py
python3 -B code/endpoint.py coefficients --p 2 --order 4
python3 -B code/endpoint.py coefficients --p 3 --order 6
python3 -B code/endpoint.py critical --order 4
python3 -B code/endpoint.py marked --p 3 --order 3 --cumulant 2
\end{verbatim}
The core never infers coefficients by fitting numerical values.
The rational-polynomial results can also be requested as JSON.
The ordinary selftest itself launches a separate optimized-Python guard
check; invoking the whole test under \texttt{-O} additionally checks that
its correctness tests do not depend on removable assertions.

With the optional dependency in \texttt{code/requirements-optional.txt},
run
\begin{verbatim}
python3 -B code/selftest.py --numerical
python3 -B code/diagnostics.py suite --deep
\end{verbatim}
The deep suite reproduces the complete $n=10000$ sums and the marked
$n=10^9$ diagnostics, together with finite Newton experiments.
Its JSON explicitly labels all transcendental values uncertified and
records the $p=3$ preasymptotic warning.

To verify the frozen file inventory and rebuild the document, use
\begin{verbatim}
python3 -B build.py --verify-only
python3 -B build.py --output ../rebuilt
python3 -B -O build.py --output ../rebuilt-optimized --numerical
\end{verbatim}
A PDF build needs \texttt{pdflatex} with the standard packages listed in
\texttt{article.tex}. The output directory must lie outside the extracted
package, preserving its frozen inventory. The build runs exact checks,
compiles in a disposable directory, and writes a PDF, deterministic ZIP,
and check summary to the chosen output directory. With \texttt{--numerical},
it also reproduces and compares the supplied high-precision diagnostics.

Deterministic timestamps and archive metadata are used. Byte-identical
PDF rebuilding is checked on the documented TeX toolchain; another TeX
release can produce different PDF bytes even when the mathematical content
is identical. The source manifest still verifies the original frozen
files independently. The build distinguishes content verification from
claims about universal cross-toolchain PDF identity.

\subsection*{Useful exact recurrences}
The public code uses only elementary rational polynomial operations.
For the power sums $S_j(r)=\sum_{h=0}^{r-1}h^j$, telescoping
$r^{j+1}=\sum_{h=0}^{r-1}((h+1)^{j+1}-h^{j+1})$ gives
\[
 S_j(r)=\frac1{j+1}\left(r^{j+1}
                -\sum_{i=0}^{j-1}\binom{j+1}{i}S_i(r)\right).
\]
The Touchard recurrence is
$T_{d+1}(z)=z(T_d(z)+T_d'(z))$, with $T_0=1$.
Together with \eqref{spb:ep:eq:Dgeneral}, \eqref{spb:ep:eq:Hrec}, \eqref{spb:ep:eq:Lrec}, and
\eqref{spb:ep:eq:assocStirling}, these determine every displayed coefficient.

The public implementation caps resources to keep accidental requests from
creating unbounded computations. Extending those caps is a programming
choice; it is not evidence of an asymptotic theorem with growing order.
\begin{writenote}[The shipped files, 7 October 2026]
In ProveIt the PDF and \file{MANIFEST.sha256} are not shipped; the programs
are \file{code/230-endpoint-*.py} and their recorded outputs and requirements
are \file{data/230-endpoint-*}. The programs import each other under their
delivered names, so the report README gives commands that rerun them in a
scratch copy with the prefix removed; the builder expects the delivered layout
and runs only in a re-extracted archive.
\end{writenote}

\begin{thebibliography}{9}
\bibitem{oeis1} The OEIS Foundation Inc., \emph{A360592},
\url{https://oeis.org/A360592}. Indexed primary-entry formula observed 5 October 2026;
see Section~\ref{sec:sources} for retrieval limitations.
\bibitem{oeis2} The OEIS Foundation Inc., \emph{A360479},
\url{https://oeis.org/A360479}. Primary-entry formula observed 5 October 2026;
see Section~\ref{sec:sources}.
\bibitem{oeis3} The OEIS Foundation Inc., \emph{A360747},
\url{https://oeis.org/A360747}. Indexed primary-entry formula observed 5 October 2026;
see Section~\ref{sec:sources}.
\bibitem{dlmf} NIST Digital Library of Mathematical Functions,
\emph{Lagrange Inversion Theorem}, Section 1.10(vii),
\url{https://dlmf.nist.gov/1.10.vii};
\emph{Charlier Polynomials}, equation 18.23.5,
\url{https://dlmf.nist.gov/18.23.E5}.
\bibitem{corless} R. M. Corless, G. H. Gonnet, D. E. G. Hare,
D. J. Jeffrey, and D. E. Knuth, \emph{On the Lambert $W$ Function},
Advances in Computational Mathematics \textbf{5} (1996), 329--359,
\url{https://doi.org/10.1007/BF02124750}.
\end{thebibliography}
\end{document}
